\documentclass[11pt]{article}
\usepackage[a4paper,margin=22mm]{geometry}
\usepackage{fontspec}
\setmainfont{DejaVu Serif}
\setsansfont{DejaVu Sans}
\usepackage{microtype}
\usepackage{xcolor}
\usepackage{hyperref}
\usepackage{enumitem}
\usepackage{parskip}
\definecolor{Accent}{HTML}{971D32}
\definecolor{Muted}{HTML}{666666}
\hypersetup{colorlinks=true,urlcolor=Accent,linkcolor=Accent}
\setlist{leftmargin=1.6em,itemsep=0.3em,topsep=0.35em}
\setcounter{tocdepth}{2}
\renewcommand{\contentsname}{Contents}
\newcommand{\sectionrule}{\par\vspace{-0.5em}\noindent\textcolor{Accent}{\rule{\linewidth}{0.6pt}}\par}
\begin{document}

\begin{center}
{\sffamily\bfseries\color{Accent} TOEFL WORD OF THE DAY\par}
\vspace{8mm}
{\fontsize{42}{48}\selectfont\bfseries connotation\par}
\vspace{2mm}
{\Large\sffamily\color{Accent} /ˌkɑːnəˈteɪʃən/\par}
\vspace{2mm}
{\small\sffamily Local audio phrase: connotation\par}
{\small\sffamily Audio file: audios/connotation.mp3\par}
\vspace{2mm}
{\large\sffamily\color{Muted} countable noun\par}
\vspace{5mm}
{\sffamily an additional association of a word \textbullet\ an implied idea or attitude\par}
\vspace{6mm}
{\large A connotation is an additional idea, feeling, or attitude suggested by a word, object, action, or situation.\par}
\vspace{6mm}
\begin{minipage}{0.84\textwidth}
\itshape “Although slim and skinny have similar meanings, skinny often has a more negative connotation.”
\end{minipage}\par
\vspace{10mm}
\textbf{Date:} August 4th 2026\quad\textbullet\quad \textbf{Level:} B1--mid-B2 TOEFL
\vfill
Created and compiled by \textbf{Amir Kooshky}\par
\vspace{2mm}
\href{https://t.me/KooshkyTOEFL}{Telegram Channel}\quad\textbullet\quad
\href{https://www.instagram.com/kooshkytoefl}{Instagram}\quad\textbullet\quad
\href{https://t.me/KooshkyTOEFL\_pv}{Personal Telegram}
\end{center}
\newpage
\tableofcontents
\newpage

\section{Meanings and usage}
\sectionrule
\subsection{Meaning 1: Additional association of a word}
\textbf{Part of speech:} countable noun

\textbf{IPA:} /ˌkɑːnəˈteɪʃən/

\textbf{Pronunciation phrase:} connotation

\textbf{Local audio file:} audios/connotation.mp3

\textbf{Register:} neutral to formal; common in linguistics, literature, communication, cultural studies, and academic writing

\textbf{Definition:} an idea or feeling that a word suggests in addition to its direct meaning

\textbf{In simpler words:} an extra feeling or idea connected to a word

\textbf{Typical pattern:} a positive, negative, or emotional connotation

\subsubsection{Natural examples}
\begin{enumerate}
  \item The word home often has connotations of comfort, safety, and belonging.
  \item Although slim and skinny have similar meanings, skinny often has a more negative connotation.
  \item In many cultures, the color white has connotations of purity and peace.
  \item The term reform usually has a positive connotation because it suggests improvement.
  \item The word childish has a more negative connotation than childlike.
  \item Translators must consider both the meaning and the connotations of a word.
  \item The expression carries emotional connotations that are difficult to reproduce in another language.
\end{enumerate}

\subsubsection{Nuance and implication}
\textbf{Central nuance:} A connotation is not usually part of a word's basic dictionary definition. It is an additional association that people connect with the word.

\textbf{Usual implication:} Connotations may be positive, negative, neutral, emotional, social, or cultural, and they can vary between groups and contexts.

\textbf{Compared with denotation:} Denotation is the direct or basic meaning of a word. Connotation is an additional idea, feeling, or association suggested by that word.

\subsubsection{Grammar notes}
\textbf{How to use it:} Use a connotation for one particular association and connotations for several associations.

\textbf{What can follow it:} It is often followed by of and an idea or feeling, as in connotations of wealth and power.

\textbf{Common preposition:} Use the connotation of a word when discussing the word's association. Use connotations of something to name the ideas it suggests.

\textbf{Noun use:} This is usually a countable noun. The plural connotations is very common because a word may suggest several different ideas or feelings.

\subsubsection{Useful patterns}
\textbf{Pattern:} have a positive or negative connotation

\textbf{Use:} Use this to describe whether the feeling associated with a word is favorable or unfavorable.

\textbf{Example:} The word ambitious usually has a positive connotation in professional contexts.

\textbf{Pattern:} carry a connotation of + noun

\textbf{Use:} Use this to identify an additional idea or feeling suggested by a word or expression.

\textbf{Example:} The phrase carries a connotation of social superiority.

\textbf{Pattern:} the connotations of + word or expression

\textbf{Use:} Use this when discussing the associations connected with a particular term.

\textbf{Example:} The article examines the political connotations of the word freedom.

\textbf{Pattern:} connotations of + idea or feeling

\textbf{Use:} Use this to name the particular associations suggested.

\textbf{Example:} The image has connotations of isolation and loss.

\subsubsection{Common wrong usage}
\paragraph{Correction 1}
\textbf{Wrong:} The word inexpensive has a good connotation than cheap.

\textbf{Correct:} The word inexpensive has a more positive connotation than cheap.

\textbf{Why:} Use more positive or more negative when comparing connotations.

\paragraph{Correction 2}
\textbf{Wrong:} The connotation of snake is a long, legless reptile.

\textbf{Correct:} The denotation of snake is a long, legless reptile.

\textbf{Why:} A direct dictionary meaning is a denotation. A connotation would be an associated idea such as danger or dishonesty.

\paragraph{Correction 3}
\textbf{Wrong:} This word connotates danger.

\textbf{Correct:} This word connotes danger.

\textbf{Why:} The verb is connote, not connotate.

\subsection{Meaning 2: Implied idea or attitude}
\textbf{Part of speech:} countable noun

\textbf{IPA:} /ˌkɑːnəˈteɪʃən/

\textbf{Pronunciation phrase:} connotation

\textbf{Local audio file:} audios/connotation.mp3

\textbf{Register:} neutral to formal; common in social, political, historical, artistic, and cultural analysis

\textbf{Definition:} an additional idea, attitude, or implication associated with an object, action, situation, or form of behavior

\textbf{In simpler words:} an extra idea or message suggested by something

\textbf{Typical pattern:} have or carry social, political, cultural, or religious connotations

\subsubsection{Natural examples}
\begin{enumerate}
  \item For some employees, working from home has connotations of freedom and flexibility.
  \item The building's grand design carries connotations of power and authority.
  \item Refusing the invitation might have unintended social connotations.
  \item The uniform has strong military connotations.
  \item The decision to close the meeting had political connotations that were not immediately obvious.
  \item In that society, owning land had connotations of wealth and social status.
  \item The gesture has religious connotations in some communities.
\end{enumerate}

\subsubsection{Nuance and implication}
\textbf{Central nuance:} In this broader use, the association comes from more than a word. An action, object, image, style, or situation can also suggest an additional message.

\textbf{Usual implication:} The association may not have been intended, but people can still interpret or respond to it.

\textbf{Compared with consequence:} A consequence is a result that happens because of an action or event. A connotation is an idea or attitude that the action or event suggests.

\subsubsection{Grammar notes}
\textbf{How to use it:} Use it to describe the additional ideas or attitudes associated with an action, object, image, or situation.

\textbf{What can follow it:} It often follows adjectives such as social, political, cultural, emotional, historical, religious, or sexual.

\textbf{Common preposition:} Use connotations of to name the ideas suggested, as in connotations of authority. Use for when identifying the group that makes the association, as in negative connotations for local residents.

\textbf{Noun use:} The plural connotations is especially common in this broader meaning because something may suggest several related ideas.

\subsubsection{Useful patterns}
\textbf{Pattern:} have social, political, or cultural connotations

\textbf{Use:} Use this when something suggests ideas connected with society, politics, or culture.

\textbf{Example:} The clothing style has political connotations in that region.

\textbf{Pattern:} carry connotations of + noun

\textbf{Use:} Use this to name the additional ideas communicated by an object, action, or design.

\textbf{Example:} The ceremony carries connotations of unity and national identity.

\textbf{Pattern:} unintended connotations

\textbf{Use:} Use this when something suggests an idea that the speaker, writer, or designer did not mean to communicate.

\textbf{Example:} The advertisement had unintended connotations of social exclusion.

\textbf{Pattern:} be free of negative connotations

\textbf{Use:} Use this when something does not carry unfavorable associations.

\textbf{Example:} The researchers wanted a neutral label that was free of negative connotations.

\subsubsection{Common wrong usage}
\paragraph{Correction 1}
\textbf{Wrong:} The gesture has a political consequence of protest.

\textbf{Correct:} The gesture has political connotations of protest.

\textbf{Why:} Use connotation for an idea suggested by the gesture. Consequence refers to an actual result.

\paragraph{Correction 2}
\textbf{Wrong:} The design is connotation with wealth.

\textbf{Correct:} The design has connotations of wealth.

\textbf{Why:} Use have connotations of, not be connotation with.

\paragraph{Correction 3}
\textbf{Wrong:} His clothing carried a connotation about authority.

\textbf{Correct:} His clothing carried a connotation of authority.

\textbf{Why:} Use a connotation of something.

\section{Word family}
\sectionrule
\subsection{connote}
\textbf{Part of speech:} transitive verb

\textbf{IPA:} /kəˈnoʊt/

\textbf{Definition:} to suggest an additional idea, feeling, or association beyond a direct meaning

\textbf{Relationship to the headword:} This verb describes suggesting the additional association called a connotation.

\textbf{Grammar pattern:} word, image, action, or object + connotes + idea or feeling

\textbf{Usage note:} Do not confuse connote with denote, which refers to direct meaning or representation.

\subsubsection{Examples}
\begin{enumerate}
  \item The word luxury often connotes wealth and exclusivity.
  \item In some contexts, silence may connote agreement.
\end{enumerate}

\subsection{connotative}
\textbf{Part of speech:} adjective

\textbf{IPA:} /kəˈnoʊtətɪv/

\textbf{Definition:} relating to the additional ideas or feelings suggested by a word or sign

\textbf{Relationship to the headword:} This adjective describes meaning or language based on connotations.

\textbf{Grammar pattern:} connotative + meaning, language, association, or use

\textbf{Usage note:} It is mainly used in linguistics, language teaching, and literary analysis.

\subsubsection{Examples}
\begin{enumerate}
  \item The lesson compared the literal and connotative meanings of several words.
  \item Poetry often depends heavily on connotative language.
\end{enumerate}

\section{Common collocations}
\sectionrule
\subsection{positive connotation}
\textbf{Applies to:} Additional association of a word

\textbf{Meaning:} a favorable idea or feeling associated with a word

\textbf{Pattern:} have a positive connotation

\textbf{Example:} The word innovative usually has a positive connotation.

\subsection{negative connotation}
\textbf{Applies to:} Additional association of a word

\textbf{Meaning:} an unfavorable idea or feeling associated with a word

\textbf{Pattern:} have a negative connotation

\textbf{Example:} The term propaganda often has a negative connotation.

\subsection{emotional connotation}
\textbf{Applies to:} Additional association of a word

\textbf{Meaning:} a feeling suggested by a word in addition to its direct meaning

\textbf{Example:} The word childhood has a strong emotional connotation for many people.

\subsection{cultural connotations}
\textbf{Applies to:} all senses

\textbf{Meaning:} associations created by the beliefs and traditions of a culture

\textbf{Example:} Colors may have different cultural connotations in different societies.

\subsection{political connotations}
\textbf{Applies to:} Implied idea or attitude

\textbf{Meaning:} associations connected with political ideas, groups, or conflicts

\textbf{Example:} The symbol acquired political connotations during the protest movement.

\subsection{social connotations}
\textbf{Applies to:} Implied idea or attitude

\textbf{Meaning:} associations connected with social identity, status, or behavior

\textbf{Example:} The choice of clothing may carry social connotations related to class or profession.

\subsection{religious connotations}
\textbf{Applies to:} Implied idea or attitude

\textbf{Meaning:} associations connected with religious beliefs or practices

\textbf{Example:} The image has religious connotations in several cultures.

\subsection{carry connotations of}
\textbf{Applies to:} all senses

\textbf{Meaning:} suggest additional ideas or feelings

\textbf{Pattern:} carry connotations of + idea or feeling

\textbf{Example:} The phrase carries connotations of authority and control.

\subsection{have a connotation of}
\textbf{Applies to:} all senses

\textbf{Meaning:} be associated with an additional idea or feeling

\textbf{Pattern:} have a connotation of + noun

\textbf{Example:} The term has a connotation of unnecessary complexity.

\subsection{unintended connotations}
\textbf{Applies to:} Implied idea or attitude

\textbf{Meaning:} associations that were not deliberately communicated

\textbf{Example:} The slogan had unintended connotations that offended some readers.

\subsection{the connotations of a word}
\textbf{Applies to:} Additional association of a word

\textbf{Meaning:} the additional ideas and feelings associated with a word

\textbf{Pattern:} the connotations of + word or expression

\textbf{Example:} Language learners should study the connotations of a word before using it in formal writing.

\section{Synonyms and antonyms}
\sectionrule
\subsection{Close synonyms}
\subsubsection{association}
\textbf{Part of speech:} countable noun

\textbf{Applies to:} all senses

\textbf{Shared meaning:} Both words can describe an idea or feeling mentally connected with something.

\textbf{Difference:} Association is a broad word for any mental connection between things. Connotation specifically refers to an additional idea or feeling suggested by a word, sign, object, or action.

\textbf{Example:} The smell of the ocean has strong associations with childhood for her.

\subsubsection{implication}
\textbf{Part of speech:} countable noun

\textbf{Applies to:} Implied idea or attitude

\textbf{Shared meaning:} Both can refer to an additional message that is not expressed directly.

\textbf{Difference:} An implication is a meaning or conclusion suggested without being stated directly. A connotation is an association or feeling connected with a word, object, or action.

\textbf{Example:} His refusal had the implication that he did not support the proposal.

\subsubsection{undertone}
\textbf{Part of speech:} countable noun

\textbf{Applies to:} Implied idea or attitude

\textbf{Shared meaning:} Both can describe an indirect attitude or feeling.

\textbf{Difference:} An undertone is a subtle attitude or quality beneath the main message, often in speech, writing, or behavior. A connotation is a more general association attached to a word or other sign.

\textbf{Example:} Her remarks contained an undertone of criticism.

\section{Words learners may confuse}
\sectionrule
\subsection{denotation}
\textbf{Part of speech:} countable or uncountable noun

\textbf{Why learners confuse them:} The two terms are commonly taught together because they describe different aspects of meaning.

\textbf{Key difference:} Denotation is the direct or basic meaning of a word. Connotation is an additional idea, emotion, or association suggested by the word.

\textbf{connotation:} The word rose may have connotations of love and romance.

\textbf{denotation:} The denotation of rose is a type of flowering plant.

\subsection{implication}
\textbf{Part of speech:} countable noun

\textbf{Why learners confuse them:} Both words describe ideas communicated indirectly.

\textbf{Key difference:} A connotation is an established or common association connected with a word or sign. An implication is a conclusion or message suggested by a particular statement, action, or situation.

\textbf{connotation:} The word cheap often has a negative connotation.

\textbf{implication:} His sudden resignation had the implication that serious problems existed.

\section{Etymology}
\sectionrule
\textbf{Origin:} Connotation comes from Latin connotare, meaning to mark or note in addition.

\textbf{Development:} The word developed the sense of an additional meaning or association attached to a word beyond its direct meaning.

\section{Practice exercises}
\sectionrule
\subsection{Word family choice}
Choose the best member of the word family for each blank.

\begin{enumerate}
  \item The word elegant often \_\_\_\_\_ beauty, style, and good taste.
  \item Poets often choose words for both their direct and \_\_\_\_\_ meanings.
  \item The term has a negative \_\_\_\_\_ in many professional settings.
\end{enumerate}

\subsection{Correct or incorrect?}
Decide whether each sentence is correct. Correct the inaccurate ones.

\begin{enumerate}
  \item The word confident usually has a more positive connotation than arrogant.
  \item The connotation of the word tree is a tall plant with a wooden trunk.
  \item The image carries cultural connotations of renewal and hope.
  \item The word connotates social status.
\end{enumerate}

\subsection{Collocation practice}
Complete each sentence with a natural collocation.

\begin{enumerate}
  \item The word stubborn often has a negative \_\_\_\_\_.
  \item The color gold carries connotations \_\_\_\_\_ wealth and success.
  \item The symbol has strong political \_\_\_\_\_ in the region.
  \item Before using the expression, consider whether it has any unintended \_\_\_\_\_.
\end{enumerate}

\subsection{Complete answer key}
\subsubsection{Word family choice}
\begin{enumerate}
  \item \textbf{connotes.} The word elegant often connotes beauty, style, and good taste. \emph{Use the verb connote because the word suggests these additional ideas.}
  \item \textbf{connotative.} Poets often choose words for both their direct and connotative meanings. \emph{Use the adjective connotative before meanings.}
  \item \textbf{connotation.} The term has a negative connotation in many professional settings. \emph{Use the noun connotation after a negative.}
\end{enumerate}

\subsubsection{Correct or incorrect?}
\begin{enumerate}
  \item \textbf{Correct} \emph{The sentence correctly compares the associations of two words.}
  \item \textbf{Incorrect}. The denotation of the word tree is a tall plant with a wooden trunk. \emph{A direct dictionary meaning is a denotation, not a connotation.}
  \item \textbf{Correct} \emph{Carry connotations of is a natural pattern for additional associations.}
  \item \textbf{Incorrect}. The word connotes social status. \emph{The correct verb is connote.}
\end{enumerate}

\subsubsection{Collocation practice}
\begin{enumerate}
  \item \textbf{connotation.} The word stubborn often has a negative connotation. \emph{A negative connotation is an unfavorable association.}
  \item \textbf{of.} The color gold carries connotations of wealth and success. \emph{Use carry connotations of something.}
  \item \textbf{connotations.} The symbol has strong political connotations in the region. \emph{Political connotations are associations connected with political ideas or groups.}
  \item \textbf{connotations.} Before using the expression, consider whether it has any unintended connotations. \emph{Unintended connotations are associations that the speaker did not mean to communicate.}
\end{enumerate}

\vfill
\begin{center}
\small\color{Muted} Created and compiled by \textbf{Amir Kooshky}\\
\href{https://t.me/KooshkyTOEFL}{Telegram Channel}\quad\textbullet\quad
\href{https://www.instagram.com/kooshkytoefl}{Instagram}\quad\textbullet\quad
\href{https://t.me/KooshkyTOEFL\_pv}{Personal Telegram}
\end{center}
\end{document}
